%----------------------------------------------------------------------------------------
% Tailored CV — Deep Learning Engineer, Omics Data / Cancer & Aging Genomics
% Poetsch Group, Institut für Genetik / CECAD, Universität zu Köln
% Ref. Wiss2604-09
%----------------------------------------------------------------------------------------

\documentclass[10pt,a4paper,sans]{moderncv}

\usepackage[utf8]{inputenc}
\usepackage[T1]{fontenc}

\moderncvstyle{classic}
\moderncvcolor{blue}

\usepackage[scale=0.78]{geometry}
\usepackage{ragged2e}

%----------------------------------------------------------------------------------------

\firstname{Kundai Farai}
\familyname{Sachikonye}

\address{Auf der Lichtung 19}{82178 Puchheim, Germany}
\mobile{(+49) 1626386344}
\email{kundai.sachikonye@bitspark.com}
\homepage{github.com/fullscreen-triangle}

%----------------------------------------------------------------------------------------

\begin{document}

\makecvtitle

%----------------------------------------------------------------------------------------
%	EDUCATION
%----------------------------------------------------------------------------------------

\section{Education}

\cventry{2019--2021}{PhD candidate, Computational Lipidomics}{Technische Universität München / Bayerisches Forschungsinstitut}{}{}{
  ML for large-scale omics data, multi-omics statistical modelling, distributed
  deep learning for clinical datasets. Departed to pursue independent research.
}
\cventry{2018--2019}{MSc Computational Biology}{Universität Tübingen}{}{}{
  Deep learning for biological sequences, statistical bioinformatics,
  genomics pipeline development.
}
\cventry{2014--2017}{MSc Pharmaceutical Biotechnology}{HAW Hamburg}{}{}{}
\cventry{2011--2014}{BSc Biochemistry and Cell Biology}{Jacobs University Bremen}{}{}{}

%----------------------------------------------------------------------------------------
%	RELEVANT MANUSCRIPTS
%----------------------------------------------------------------------------------------

\section{Relevant Manuscripts}

\cvitem{2025}{
  \textbf{On the Thermodynamic Consequences of Categorical Completion on Biological Systems.}
  CNN/DFT pipeline for MHC binding prediction: AUC 0.81 vs NetMHCpan 0.68
  ($r=0.73$, $p<10^{-12}$, $n=2{,}847$ IEDB peptides).
  \textit{Manuscript.}
}
\cvitem{2025}{
  \textbf{On Disease Epistemology in Blind Receiver.}
  Deep learning diagnostic framework; AUC$=1.00$; 25/25 validation experiments;
  sparse LP for personalised therapeutic target selection.
  \textit{Manuscript. AIMe Registry / TUM Weihenstephan.}
}

%----------------------------------------------------------------------------------------
%	DEEP LEARNING AND OMICS TOOLS
%----------------------------------------------------------------------------------------

\section{Deep Learning and Omics Tools}

\cvitem{gospel}{
  \textbf{Genomics Deep Learning Pipeline.}
  $10^6$ variants/second VCF processing; somatic variant calling for cancer
  genomics; RNA-seq quantification and differential expression; Bayesian network
  inference and GMMs for multi-omics integration; pathway analysis (KEGG,
  Reactome, GSEA). Validated: 1000~Genomes Phase~3, gnomAD~v3.1.2, UK~Biobank.
  ClinVar / COSMIC API integration. Ray distributed, Docker deployable.
  \href{https://github.com/fullscreen-triangle/gospel}{github.com/fullscreen-triangle/gospel}
}

\cvitem{lavoisier}{
  \textbf{CNN + Classical Dual-Pipeline Omics Analysis.}
  Dual pipeline: numerical MS1/MS2 processing (peak detection, spectral library
  matching) alongside CNN visual pipeline (PyTorch) converting spectra to
  temporal image sequences. 98.9\,\% NIST accuracy; Ray distributed; $>$100\,GB
  mzML support; memory-mapped I/O.
  Demonstrates paper-to-production DL implementation: academic spectral imaging
  literature $\to$ deployed genomics-scale analysis tool.
  Virtual instrument: \href{https://lavoisier.vercel.app/experiment}{lavoisier.vercel.app}.
  \href{https://github.com/fullscreen-triangle/lavoisier}{github.com/fullscreen-triangle/lavoisier}
}

\cvitem{syndrome}{
  \textbf{Deep Learning Toolkit for Cancer Immunopeptidology.}
  Live annotation tools at \href{https://syndrome-seven.vercel.app/tools}{syndrome-seven.vercel.app}:
  MHC binding predictor (AUC 0.81 vs NetMHCpan 0.68);
  neoantigen identifier (dual threshold: $\Delta\rho_\mathrm{self}>0.15$ + $\rho_\mathrm{MHC}>0.72$);
  immune escape predictor (systematic single-aa substitution scan);
  antibiotic resistance predictor (drug-enzyme spectral complementarity).
  $\mathcal{O}(cL\log L)$; browser-deployable annotation tools for lab use.
  \href{https://github.com/fullscreen-triangle/syndrome}{github.com/fullscreen-triangle/syndrome}
}

\cvitem{four-sided-triangle}{
  \textbf{Automated Evaluation Framework + Specialised DL Library.}
  8-stage metacognitive pipeline with automated benchmarking at each stage;
  Rust core specialised library (15--50$\times$ speedup, PyO3 bindings);
  FastAPI REST interface; Docker/Kubernetes; GQIC resource allocation.
  Template for building automated evaluation frameworks and specialised
  DL libraries for the lab's genomics infrastructure.
  \href{https://github.com/fullscreen-triangle/four-sided-triangle}{github.com/fullscreen-triangle/four-sided-triangle}
}

\cvitem{olduvai}{
  \textbf{Aging Biology: Biological Motion Deep Learning.}
  Deep learning for aging phenotype classification: AUC$=1.00$ for Parkinson's /
  ataxia / healthy ageing from wearable data (86 nights, $n$ = validated cohort).
  Circadian and motor oscillator dimensions directly connected to genomic
  expression changes in aging. Relevant to Poetsch group's aging genomics programme.
  \href{https://github.com/fullscreen-triangle/olduvai}{github.com/fullscreen-triangle/olduvai}
}

\cvitem{purpose}{
  \textbf{Backward Training Principle --- DL Model Training Framework.}
  Formal proof: training at the hardest (pure-intent) instance propagates
  $\epsilon$-competence to all constrained sub-tasks by feasibility inclusion.
  $(N{+}1)\times$ sample savings; $2^{N-1}\times$ for binary cascades.
  Applicable to training cancer/aging genomics DL models: train at maximum
  variant density / clearest disease signal; competence propagates to
  sparse-data and noisy clinical settings.
  \href{https://github.com/fullscreen-triangle/purpose}{github.com/fullscreen-triangle/purpose}
}

%----------------------------------------------------------------------------------------
%	EXPERIENCE
%----------------------------------------------------------------------------------------

\section{Experience}

\cventry{2021--present}{Independent AI \& Computational Biology Developer}{Bitspark GmbH}{}{}{
  Deep learning systems for genomics, transcriptomics, immunopeptidology, and
  aging biology. Algorithm implementation from academic papers to production
  deployment. Internal tools, automated evaluation frameworks, specialised
  libraries. All systems publicly deployed and validated.
}

\cventry{2019--2021}{Computational Lipidomics}{TUM / Bayerisches Forschungsinstitut}{Munich}{}{
  ML for large-scale MS omics data: CNN spectral classification, distributed
  pipeline maintenance, multi-omics integration for clinical cohorts. Python, Java.
}

\cventry{2017--2018}{Glycomics and Proteomics}{Universitätsklinikum Hamburg-Eppendorf}{}{}{
  Automated LC-MS/MS and MALDI-TOF pipeline; SOP development; clinical sample
  processing; data annotation and quality control.
}

\cventry{2013}{Cancer Biomarker Discovery}{University Medical Center Groningen}{}{}{
  UPLC-MS/MS shotgun proteomics for cancer biomarker discovery; ML-based
  statistical validation.
}

%----------------------------------------------------------------------------------------
%	TECHNICAL SKILLS
%----------------------------------------------------------------------------------------

\section{Technical Skills}

\cvitem{Deep Learning / ML}{
  PyTorch (primary); CNNs, LSTMs, Transformers (1D and 2D), GNNs;
  knowledge distillation; variational Bayes; Bayesian networks; GMMs;
  Gaussian processes; surrogate model training; RAG pipelines;
  LoRA / parameter-efficient fine-tuning
}

\cvitem{Genomics / Omics DL}{
  Variant effect prediction; somatic variant calling; RNA-seq DL quantification;
  spectral deep learning for MS data; sequence-to-function models;
  neoantigen / MHC binding prediction; multi-omics integration
}

\cvitem{Internal Tools / Libraries}{
  Rust specialised libraries with PyO3 Python bindings (15--50$\times$ speedup);
  automated evaluation frameworks with reference database benchmarking;
  browser-deployable annotation tools (WebGPU, Vercel);
  FastAPI REST interfaces; Docker / Kubernetes deployment
}

\cvitem{Data / Infrastructure}{
  Python, R, Rust, Java, TypeScript, Bash;
  Ray, Dask; HDF5, mzML, VCF, FASTQ, BAM;
  memory-mapped I/O for $>$100\,GB datasets;
  gnomAD, ClinVar, COSMIC, KEGG, Reactome API clients; Git, CI/CD
}

%----------------------------------------------------------------------------------------
%	LANGUAGES
%----------------------------------------------------------------------------------------

\section{Languages}

\cvitemwithcomment{English}{Native / Fluent}{}
\cvitemwithcomment{German}{Conversational}{Living in Germany since 2011}

%----------------------------------------------------------------------------------------

\renewcommand{\listitemsymbol}{-~}

\end{document}
